\documentclass[10pt,a4paper]{amsart}
\usepackage[margin=19mm]{geometry}
\usepackage{amsmath,amssymb,mathtools}
\usepackage[scr=rsfs,cal=euler]{mathalfa}
\usepackage{xfrac}
\usepackage[unicode,hidelinks,bookmarks=false]{hyperref}
\newcommand{\ie}{i.e.\ }
\newcommand{\cf}{cf.\ }
\newcommand{\resp}{respectively\ }
\newcommand{\Subsection}{\subsection}
\providecommand{\nobreakdash}{\nobreak\hbox{-}}
\setlength{\emergencystretch}{2em}
\multlinegap0pt
\usepackage{bm}
\usepackage{bbm}
\usepackage{dsfont}
\usepackage{xspace}
\usepackage{cancel}
\usepackage{mathtools}
\usepackage{amsthm}
\makeatletter
\@ifundefined{thm}{\newtheorem{thm}{Thm}}{}
\@ifundefined{prop}{\newtheorem{prop}[thm]{Proposition}}{}
\makeatother
\def\cK{{\mathcal K}}
\newcommand{\ddc}{\sqrt{-1}\partial\bar\partial}
\title[On the uniqueness of even $L^p$ Minkowski problem]{On the uniqueness of even $L^p$ Minkowski problem}
\author{Weiyong He, Junbang Liu}
\date{Source: arXiv:2510.21530v1 (CC BY 4.0)}
\begin{document}
\maketitle

\noindent\emph{Source and licence.} On the uniqueness of even $L^p$ Minkowski problem. Version arXiv:2510.21530v1 (CC BY 4.0), distributed under the Creative Commons Attribution 4.0 International licence, as recorded in the arXiv licence metadata for this exact version and shown on the abstract page https://arxiv.org/abs/2510.21530v1. Full rights notice: licenses/arxiv-2510.21530-CC-BY-4.0.txt.\par\smallskip

\noindent\textit{Editorial scope.} Counted unit: the Prop whose source label is \texttt{None} in the article (source0.tex lines 921--923, source label \texttt{None}), delivered together with the article's own complete original proof of it (source0.tex lines 924--995, one \texttt{proof} environment of 72 lines). No mathematics is written, repaired or re-derived: statement and proof are the article's own text line for line. Required context retained uncounted: none. Every definition, display and label of those lines is retained unchanged. Omitted from this package and not redistributed: 1--920, 996--1234. Nothing inside a retained excerpt is omitted or altered: every retained body is delivered exactly as the article's lines are written, so no omission receipt of any kind accompanies this package. Citations: the retained excerpts carry no citation command at all, so no bibliography entry of the article is transcribed and no reference is redistributed. Reference adaptation: none was needed and none was made; every reference inside the delivered unit resolves inside it, so no unresolved reference or citation is produced. Printed slips kept literal: no printed word, symbol, hypothesis or number of the counted statement or of its proof was altered. Layout and class: the article's own class is \texttt{amsart} with packages including color, amsthm, amsmath, amssymb, amscd, graphics, latexsym, stmaryrd, hyperref, empheq; the wrapper is the lane's pinned \texttt{amsart} editorial wrapper at 10pt on A4, which restates the theorem-like environments the retained text needs and adds the article's own macro definitions that the retained text needs (\texttt{\textbackslash cK}, \texttt{\textbackslash ddc}), with extra wrapper packages (bm, bbm, dsfont, xspace, cancel, mathtools). The delivered numbering is the unit's own. Assets: no retained span contains a figure environment, a graphics-inclusion command, a PDF-page inclusion, a table or a TikZ picture; no image file is copied into this package, the package declares no asset dependency and no image file is redistributed with it. Licence: version arXiv:2510.21530v1 of this article is distributed under the Creative Commons Attribution 4.0 International licence, as recorded in the arXiv licence metadata for this exact version and shown on the abstract page https://arxiv.org/abs/2510.21530v1. A full rights notice is delivered at licenses/arxiv-2510.21530-CC-BY-4.0.txt. Characters: every retained excerpt is pure ASCII, so the delivered engine, XeLaTeX with its default Latin Modern text font, reports no missing character.
\begin{prop}
    For a convex body $K\in \cK_e$ with the support function $C^{-1}\leq h\leq C$, then $\lambda_{1, e}(-L_K)=\lambda_{1, e}(K)$. 
\end{prop}

\begin{proof}
Fix $\epsilon>0$.
    Suppose $\lambda_{1, e}(-L_K)+\epsilon$ is obtained by a smooth even function $u$, such that 
    \[
    \int_{S^n}{(u-\underline u)^2} h\det(H)dS=1.
    \]
    Now fix $u$. Choose $\tilde h$ smooth, $\|\tilde h-h\|<\epsilon_{1}(\epsilon, u)$  and denote
    \[
    \tilde \lambda=\frac{\displaystyle \int_{S^n}\tilde H^{ij}\tilde h^2u_iu_j\det(\tilde H)dS}{\displaystyle\int_{S^n}{(u-\underline u)^2} \tilde h\det(\tilde H)dS}.
    \]
    We can choose $\epsilon_1$ sufficiently small such that 
    \[
   -C\epsilon\leq \lambda_{1, e}(-L_K)-\tilde \lambda \leq C\epsilon.
    \]
    Next take $\tilde{z}=\tilde h u$, then we get 
    \[
    \tilde \lambda=\frac{\displaystyle \int_{S^n}\tilde H^{ij}\tilde{z}_i\tilde{z}_j\det(\tilde H)dS-\int_{S^n}\tilde H^{ij}\delta_{ij}\tilde{z}^2 \det{\tilde H}dS+n\int_{S^n}\tilde{z}^2\tilde h^{-1}\det{\tilde H}dS}{\displaystyle\int_{S^n}\tilde{z}^2\tilde h^{-1}\det(\tilde H)dS-\left(\int_{S^n}\tilde h\det{\tilde H}dS\right)^{-1}\left(\int_{S^n}\tilde{z}\det{\tilde H}dS\right)^2}.
    \] 
Except the term $\displaystyle \int_{S^n}\tilde H^{ij}\tilde{z}_i\tilde{z}_j\det(\tilde H)dS$, all other terms in the above are approximated by the term replacing $\tilde h$ by $h$, $\tilde{z}$ by $z=hu$. For example, we have, for $\epsilon_1$ sufficiently small, 
\[
\left|\int_{S^n}\tilde H^{ij}\delta_{ij}\tilde{z}^2 \det{\tilde H}dS-\int_{S^n} H^{ij}\delta_{ij}z^2 \det{ H}dS\right|\leq C\epsilon.
\]
To see the convergence of the first term, we use the identity $\sqrt{-1}\partial z\wedge\bar{\partial}z=\frac{1}{2}(\ddc z^2-z\ddc z)$ to rewrite it as the sum of two parts:\[
\begin{aligned}
    \int_{S^n}\tilde{H}^{ij}\tilde{z}_i\tilde{z}_j\det(\tilde{H})dS=&\int_{M}f(r)\sqrt{-1}\partial r\wedge\bar{\partial}r\wedge \sqrt{-1}\partial \tilde{z}\wedge\bar{\partial}\tilde{z}(\sqrt{-1}\ddc \tilde{h})^{n-1}\\
    =&\int_{M}f(r)\sqrt{-1}\partial r\wedge\bar{\partial}r\wedge\frac{1}{2}(\ddc \tilde{z}^2-\tilde{z}\ddc \tilde{z})\wedge(\sqrt{-1}\ddc \tilde{h})^{n-1}\\
    =:&III+IV
\end{aligned}
\]
Since $\tilde{h}\rightarrow h$ in $L^\infty$, we know that $\tilde{z}\rightarrow hu$ in $L^\infty$. Here $u$ is $C^2$ function, and $h$ is a plurisubharmonic function. Since $h$ is uniformly bounded, and it's the support function of a convex body, we have $|D h|$ is also uniformly bounded. It follows that \[
\ddc(hu)=u\ddc h+\sqrt{-1}\partial u\wedge\bar{\partial}h+\sqrt{-1}\partial{h}\wedge\bar{\partial}u+h\ddc u\geq -C_0\beta,
\] for some constant $C_0$ depending on $||u||_{C^2},||h||_{L^\infty}$. It means that $hu$ is a quasiplurisubharmonic function. Now the proof of the convergence of the integral $III,IV$ is similar to the estimate of $I$ as before, and hence, we get \[
\int_{S^n}\tilde{H}^{ij}z_iz_j\det(\tilde{H})dS\rightarrow \int_{M}f(r)\sqrt{-1}\partial r\wedge\bar{\partial}r\wedge\frac{1}{2}(\ddc z^2-z\ddc z)\wedge(\sqrt{-1}\ddc h)^{n-1}.
\] Therefore, for any $C^2$-function $\hat{z}$ with $|\hat{z}-z|_{L^\infty}\leq \epsilon_2(\epsilon,u,h)$, by the Chern-Levine-Nirenberg inequality, we have \[\left|
\int_{S}H^{ij}\hat{z}_i\hat{z}_j\det HdS-\int_{M}f(r)\sqrt{-1}\partial r\wedge\bar{\partial}r\wedge\frac{1}{2}(\ddc z^2-z\ddc z)\wedge(\sqrt{-1}\ddc h)^{n-1}\right|\leq \epsilon
\] Again, all other terms in $\tilde{\lambda}$ can be approximated by replace $\tilde{z}$ by $\hat{z}$. 
It follows that \[
\tilde{\lambda}\geq \frac{\lambda_{1,e}(K)-C\epsilon}{1+C\epsilon}.
\]Hence, we get
$\lambda_{1,e}(K)\leq \lambda_{1,e}(-L_K)$. The other direction is similar, and this proves that $\lambda_{1,e}(-L_K)=\lambda_{1,e}(K).$
%Next we claim that $$\displaystyle \int_{S^n}\tilde H^{ij}z_iz_j\det(\tilde H)dS\geq \int_{S^n}H^{ij}z_iz_j\det{H}dS-C\epsilon.$$ 
%We compute, use the complex notation, that
%\[
%\int_{S^n}\tilde H^{ij}z_iz_j\det(\tilde H)dS-\int_{S^n}H^{ij}z_iz_j\det{H}dS=c_n \int_{M} (\tilde h-h)f(r)(\sqrt{-1}\p r\wedge \bar\p r)\wedge (\sqrt{-1}\p\bar\p z)^2\wedge \Theta(h, \tilde h),
%\]
%where $\Theta(h, \tilde h)$ is given by
%\[
%\Theta(h, \tilde h)=\sum_{k} (\sqrt{-1}\p\bar\p h)^{n-2-k}\wedge (\sqrt{-1}\p\bar\p \tilde h)^{k}.
%\]
%We can choose $u$ such that $1\leq u\leq C$, and $\tilde h\geq h$ pointwise. We compute
%\[
%\sqrt{-1}\p\bar\p z=u\sqrt{-1}\p\bar\p \tilde h+\sqrt{-1}\p h\wedge \bar\partial u+\sqrt{-1}\p u\wedge \bar \p h+\tilde h^2 \sqrt{-1}\p\bar\p u.
%\]
%Since we assume that $\tilde h$ is uniformly bounded, $|\p \tilde h|$ has a uniform bound as well. It follows 
%\[
%\sqrt{-1}\p\bar\p z= u\sqrt{-1}\p\bar\p \tilde h-C\omega_0\geq -C\omega_0,
%\]
%where $\omega_0$ is the standard K\"ahler on $M$. 
%Moreover, 
%\[
%(\sqrt{-1}\p\bar\p z)^2\geq -C\omega_0\wedge \omega_0.
%\]
%Hence it follows that, given $\tilde h-h\geq 0$, that we have
%\[
%\int_{M} (\tilde h-h)(\sqrt{-1}\p r\wedge \bar\p r)\wedge (\sqrt{-1}\p\bar\p z)^2\wedge \Theta(h, \tilde h)\geq -C\int_M (\tilde h-h)f(r)(\sqrt{-1}\p r\wedge \bar\p r)\wedge \omega_0^2\wedge \Theta(h, \tilde h)\geq -C\epsilon. 
%\]
%In summary, this implies that, 
%\[
%\tilde \lambda\geq \frac{\lambda_{1, e}(K)-C\epsilon}{1+C\epsilon}.
%\]
%Hence we get that $\lambda_{1, e}(K)\leq \lambda_{1, e}(-L_K)$. The other direction is similar, and this proves that $\lambda_{1, e}(-L_K)=\lambda_{1, e}(K)$. 
\end{proof}

\end{document}
